\chapter{Feature coverage and source-only interfaces}\label{ch:coverage}
\section{Feature checklist}
This checklist maps the user-facing feature families to their detailed instructions. It includes optional and currently unexposed source implementations separately. The accompanying \code{coverage.json} records source fingerprints and UI control declarations. The check script detects changed/new source files and stale generated API or CLI reference material. It is a maintenance alarm, not an automated proof that prose explains every possible runtime state. The control descriptions were reviewed against the implementation.

\begin{longtable}{p{.47\linewidth}p{.44\linewidth}}
\toprule Feature family & Documentation \\ \midrule\endhead
Open, recent files, close, save, Save As, formats, drag/drop, tabs & Chapter~\ref{ch:start}; Section~\ref{sec:display-controls}. \\
Bulk Crystal, all crystal systems, groups, site editing & Chapter~\ref{ch:builders}; Section~\ref{sec:builder-controls}. \\
Substitutional Solid Solution, composition, seed & Chapter~\ref{ch:builders}; Section~\ref{sec:builder-controls}. \\
CSL search, grain geometry, overlap removal, reduction & Chapter~\ref{ch:builders}; Section~\ref{sec:builder-controls}. \\
All seven nanocrystal cuts and Wulff family construction & Chapter~\ref{ch:builders}; Section~\ref{sec:builder-controls}. \\
Custom OBJ/STL fill, orientation, replication & Chapter~\ref{ch:builders}; Section~\ref{sec:builder-controls}. \\
Polycrystal, random/specified/partial grain orientation & Chapter~\ref{ch:builders}; Section~\ref{sec:builder-controls}. \\
Amorphous packing, density, scale, pair constraints & Chapter~\ref{ch:builders}; Section~\ref{sec:builder-controls}. \\
Optional Stacking Faults sequence and export & Chapter~\ref{ch:builders}; Section~\ref{sec:builder-controls}. \\
Selection, undo/redo, context substitution, centroid insertion & Chapter~\ref{ch:editing}; Section~\ref{sec:edit-controls}. \\
Atom/cell table, integer transform, merge & Chapter~\ref{ch:editing}; Section~\ref{sec:edit-controls}. \\
Cell Sculptor and bounded slab intersections & Chapter~\ref{ch:editing}; Section~\ref{sec:edit-controls}. \\
Dislocation character, vectors, regions, preview/replace & Chapter~\ref{ch:editing}; Section~\ref{sec:edit-controls}. \\
Interstitial detection, type filters, random/mesh/manual insertion & Chapter~\ref{ch:editing}; Section~\ref{sec:edit-controls}. \\
Camera, colouring, representations, visibility, themes & Chapters~\ref{ch:start}, \ref{ch:analysis}; Section~\ref{sec:display-controls}. \\
Structure/Atom Info, distance, angle & Chapter~\ref{ch:analysis}; Section~\ref{sec:display-controls}. \\
RDF, shell distortion, SRO & Chapter~\ref{ch:analysis}; Section~\ref{sec:display-controls}. \\
Planes, directions, Voronoi, coordination polyhedra & Chapter~\ref{ch:analysis}; Section~\ref{sec:display-controls}. \\
Atomic radii, lighting, material, colours, image export & Chapter~\ref{ch:analysis}; Section~\ref{sec:display-controls}. \\
VASP/Cube/XSF input, quantity/units, channel selection & Chapter~\ref{ch:electronic}. \\
Electronic 3D volume/surface, 2D arbitrary slices, contours & Chapter~\ref{ch:electronic}; Section~\ref{sec:display-controls}. \\
Charge difference, masks, Boolean operations, regional charge & Chapter~\ref{ch:charge}. \\
All general electronic operations and exports & Chapter~\ref{ch:reference}. \\
Python structures, I/O, desktop/notebook viewer & Chapter~\ref{ch:python}; Appendix~\ref{ch:api-signatures}. \\
Python electronic grids, models, surfaces and tables & Chapter~\ref{ch:python}; Appendix~\ref{ch:api-signatures}. \\
All ten native builder CLI modes and Python electronic CLI & Chapter~\ref{ch:python}; Appendix~\ref{ch:api-signatures}. \\
CNA, ADF, interface matching workflows & This appendix. \\
Installation, build options, troubleshooting, example regeneration & Previous appendices and \code{docs/manual/README.md}. \\
\bottomrule
\end{longtable}

\section{Common-neighbour analysis: desktop and API}
CNA is implemented in \code{src/algorithms/CommonNeighbourAnalysis.h} and its companion source. Open Analysis > Common Neighbour Analysis, or call Python \code{af.cna}. The native \code{--analyze cna} command writes per-atom CSV results.

For C++ integration, call \code{atomforge::analysis::computeCna(structure, params)}. \code{CnaParams} supplies \code{cutoffScale} (default 1.18) and \code{usePbc} (default true). The result contains \code{valid}, \code{message}, atom/pair counts, \code{pbcUsed}, signature counts, environment counts, and per-atom rows with index, species, coordination, dominant signature/count, and environment. A signature records common-neighbour count, bonds, and chain information. Check validity before consuming these tables.

The dialog provides Use PBC when unit cell is available, Bond cutoff scale (1.00--1.60), and Run CNA. Changing cutoff changes the bond graph and therefore the classification. The fixed cutoff option is in Angstrom; zero selects the covalent-radius scale. Classification uses the complete distribution of bond signatures around each atom. Repeat small cells sufficiently to represent all neighbour sites. Dense common-neighbour graphs exceeding 16 vertices or the search budget have chain index -1 rather than an invented classification.

\section{Angular-distribution analysis: desktop and API}
ADF is implemented in \code{src/algorithms/AngularDistributionAnalysis.h}; it is available from the Analysis menu and Python \code{af.adf}. C++ callers use \code{computeADF(structure, params)} with \code{AdfParams}. The neighbour cutoff \code{rCutoff} is in \angstrom\ (default 3.5); \code{binCount} divides 0--180 degrees (default 180); \code{usePbc} controls periodic geometry. \code{centreMode} selects All, ByElement, or ByPair. Set \code{centreSymbol} for element-specific centres, plus \code{neighSymbol1} and \code{neighSymbol2} for the neighbour pair in ByPair mode.

\code{normalize} scales to peak 1, and \code{smoothPasses} sets smoothing (default 2). This normalisation is not the RDF density normalisation. \code{AdfResult} contains validity/message, echoed settings, atom/centre/triplet counts, PBC status, histogram rows (angle, raw count, displayed value), detected peaks, and per-centre-element coordination mean, standard deviation, and count.

The source dialog adds peak markers, raw-count display, plot range, and coordination statistics. Run/Compute ADF refreshes the calculation. Each unordered neighbour pair is counted once; the raw histogram remains unsmoothed even when the displayed curve is smoothed. Interpret peaks as angle distributions under the chosen neighbour definition; a label near a familiar ideal angle does not by itself identify a unique structure.

\section{Interface matching: desktop and CLI workflow}
Open Build > Interface Builder. It is distinct from Merge Structures. Its source workflow searches matching in-plane supercells of two periodic inputs, strains an interface construction, and assembles layers. The \code{--build interface} command and \code{af.build("interface", options)} support automated matching.

The source dialog loads structures A and B. Orientation relationship choices are None, Manual angle, or Plane + direction; the last takes a plane and in-plane direction for each crystal. Orientation tolerance limits acceptable angle matching. Search controls include maximum integer coefficients for each side (nmax/mmax), maximum cell counts, layers A/B, maximum mean strain, and maximum rotation. Geometry controls include Z gap, vacuum, and X/Y output repetition. The stiffness target selects A or B for a $6\times6$ stiffness matrix used by the ranking calculation. These material constants must be supplied consistently; they are not inferred from element labels.

Find Matches produces candidate combinations; inspect and select a candidate before Build. Broad coefficient/cell limits increase search work. Small reported strain does not validate chemical termination or interfacial energy. The output must be inspected for contacts and relaxed externally where appropriate.

For C++ use, \code{get2DBasis} extracts the in-plane basis; \code{generateUniqueSupercells} enumerates bounded integer supercells; \code{strainComponents} and \code{meanAbsStrain} assess matching; \code{cubicElasticDensity} evaluates the implemented elastic ranking expression. \code{makeSupercell2D}, \code{repeatLayersZ}, and \code{applyTransform2D} construct the chosen components. \code{assembleInterface} combines them with gap/vacuum, and \code{repeatInterfaceXY} repeats the product. \code{orientationAngleFromPlaneDir} derives an in-plane orientation from plane/direction data. Consult the declared types in \code{src/algorithms/InterfaceBuilder.h} when integrating; these functions are the developer interfaces underlying the desktop workflow.

\section{Keeping the documentation complete as code changes}
Run \code{python docs/manual/scripts/check_coverage.py} before building a release manual. A failure identifies a changed source fingerprint, missing source mapping, or stale API/CLI reference. Review the changed feature, update its prose, then intentionally refresh the snapshot with \code{--refresh}. Refreshing alone does not document a feature. The generated API appendix uses Python AST signatures and docstrings without importing native libraries; native CLI reference text is extracted from the help definitions and checked against the executable during this edition's validation.
